\chapter{Aspectos Conceituais} \label{cap:conceitos}

Este capítulo apresenta os fundamentos conceituais que sustentam a proposta desta pesquisa, organizando os elementos teóricos e de modelagem necessários para a compreensão do sistema. Primeiramente, descreve-se o modelo de sistema adotado, com a definição das entidades, de seus papéis institucionais e dos mecanismos criptográficos que viabilizam o controle econômico do acesso aos dados de monitoramento veicular. Em seguida, formaliza-se o modelo de adversário e as principais ameaças consideradas, destacando as hipóteses de segurança, os vetores de abuso e as estratégias de mitigação aplicáveis ao contexto de monitoramento de tráfego urbano.

\section{Modelo de Sistema}
Para endereçar os desafios de preservação da privacidade e manutenção da segurança pública no monitoramento de tráfego na cidade de São Paulo, definimos um ecossistema composto por seis entidades principais. Esta arquitetura visa distribuir a confiança entre diferentes atores municipais e estaduais para mitigar o risco de pontos únicos de falha e garantir a proporcionalidade no acesso aos dados.

O Coletor ($\mathcal{C}$) compreende os agentes de borda, representados pela rede de câmeras de monitoramento com reconhecimento automático de placas (câmeras ALPR) geridas pela Companhia de Engenharia de Tráfego (CET) ou integradas ao programa City Câmeras. O $\mathcal{C}$ é responsável pela captura inicial da placa e geração dos metadados veiculares. Antes da transmissão, o $\mathcal{C}$ realiza o "amassamento" (\textit{crumpling}) das chaves criptográficas através de funções de hash. O texto claro da placa é descartado imediatamente após a cifragem para impedir a replicação de informações pessoalmente identificáveis (PII) na origem, de modo que o sistema retenha apenas registros cuja recuperação exige esforço computacional.

O Servidor ($\mathcal{S}$) atua como o custodiante tecnológico da base de dados criptografada, função desempenhada por centros de processamento de dados como a Prodesp. O $\mathcal{S}$ armazena os registros "amassados" e gerencia as interfaces de consulta, operando estritamente no domínio cifrado. O servidor é considerado semi-honesto, executando as operações de busca sem possuir as chaves de decodificação ou a capacidade de acessar os dados de forma arbitrária sem o devido empenho de recursos.

O Usuário ($\mathcal{U}$) representa a autoridade investigativa, como as delegacias especializadas da Polícia Civil de São Paulo ou a Polícia Federal, que detém interesse legítimo em acessar registros de um veículo específico. Em investigações rotineiras, o $\mathcal{U}$ deve obrigatoriamente empenhar recursos computacionais reais, como ciclos de CPU e energia, para resolver o puzzle de prova de trabalho (proof-of-work) associado a cada registro desejado. Este mecanismo impõe um custo linear que atua como barreira técnica contra o monitoramento indiscriminado da malha viária paulistana.

O controle institucional é exercido pelo Autorizador Legal e pelo Conselho de Governança. O Autorizador Legal, tipicamente representado pelo Poder Judiciário, analisa a validade jurídica das quebras de sigilo e emite tokens de autorização necessários para legitimar o início do gasto computacional. Paralelamente, o Conselho de Governança funciona como o Auditor de Parâmetros e Integridade. Este órgão define os níveis de entropia ($l$) dos puzzles para equilibrar a viabilidade das investigações e a segurança da população, além de realizar a auditoria retroativa dos registros para verificar se o volume de decodificações processadas pelo Usuário corresponde estritamente às autorizações emitidas.

Por fim, o Comitê de Emergência é a entidade responsável pelo Acesso Excepcional em cenários de extrema gravidade. Este comitê detém fragmentos de segredos que, quando combinados através de esquemas de custódia federada, permitem resolver puzzles de abrasão de alta complexidade. Este mecanismo desbloqueia uma infraestrutura de decodificação acelerada, garantindo que o acesso em situações críticas deixe um rastro orçamentário público e auditável, impedindo suspensões sigilosas das garantias de privacidade.
    
\section{Modelo de Adversário e Ameaças}
Adotamos um modelo de ameaça semi-honesto (honest-but-curious) \cite{Link3}, no qual assume-se que o Usuário ($\mathcal{U}$) e o Servidor ($\mathcal{S}$) executam o protocolo conforme especificado, mas tentam extrair o máximo de informação possível a partir dos dados cifrados aos quais possuem acesso legítimo. A validade desta premissa no sistema de monitoramento de São Paulo sustenta-se no fato de as entidades envolvidas serem órgãos públicos sujeitos a regimes rigorosos de fiscalização e auditoria. Nesses contextos, o custo político e judicial derivado de uma má conduta detectada atua como um incentivo direto ao comportamento conforme as regras, tornando tentativas de sabotagem ativa menos prováveis que a "curiosidade" passiva sobre os dados. Além disso, a separação de papéis aliada à auditoria retroativa do Conselho de Governança assegura que qualquer tentativa de burlar as barreiras econômicas exigiria um investimento orçamentário visível e tecnologicamente proibitivo.

A segurança do protocolo fundamenta-se na Independência de Chaves, garantindo que a solução de um puzzle para um registro específico (um avistamento de placa em um determinado cruzamento) não reduza o custo de recuperação de registros subsequentes ou de outros veículos, mantendo a vigilância orçamentariamente rastreável e limitada. Assim, mesmo um adversário semi-honesto que analise exaustivamente os valores cifrados não consegue distinguir as mensagens de valores aleatórios sem o empenho do custo computacional linear previsto.

O modelo identifica e mitiga as seguintes ameaças críticas no contexto de monitoramento de tráfego na cidade de São Paulo \cite{Link4}:
\begin{itemize}
    \item \textbf{Vazamento de Informação em massa:} risco de reconstrução de trajetórias e inferência de destinos de cidadãos não-alvos. Devido à alta densidade da rede viária de São Paulo, a captura frequente permite localizar veículos com precisão preocupante (menos de 1 km de incerteza) \cite{Link5}. O sistema mitiga este risco ao converter a facilidade tecnológica de acesso em um impedimento econômico intransponível para orçamentos públicos finitos, através do custo computacional por registro.

    \item \textbf{Corrupção Institucional e Abuso de Poder:} risco de um agente isolado desviar o sistema para finalidades de vigilância política ou perseguição pessoal. A atuação do Conselho de Governança como auditor externo e a custódia de chaves para emergências garantem que nenhum indivíduo detenha o poder total de suspender as garantias de privacidade de forma sigilosa.

    \item \textbf{Adulteração e Falsificação:} tentativas de modificar registros no servidor para ocultar trajetos de suspeitos ou forjar evidências contra inocentes. Protocolos de autenticidade asseguram que apenas registros originados no Coletor ($\mathcal{C}$) e autorizados judicialmente sejam processados, protegendo a integridade da prova e evitando a inserção de dados falsos por adversários.

    \item \textbf{Ataque de Reutilização de Esforço:} ameaça de um adversário reaproveitar a solução de um registro para descriptografar outros de forma barata (ex.: ataques de dicionário ou tabelas de pré-computação). O uso de identificadores contextuais únicos, como endereços de câmeras e carimbos de tempo precisos no processo de \textit{crumpling}, garante que cada processo de decodificação seja independente, único e não-reutilizável.
\end{itemize}

O sistema é considerado seguro se o esforço necessário para realizar uma vigilância indiscriminada for tecnologicamente garantido como proibitivo, forçando a autoridade a fundamentar cada busca na necessidade investigativa e na disponibilidade orçamentária fiscalizada pelo Conselho de Governança.